 % 本文件由 scripts/assemble.py 自动装配，勿手改（改 parts/ 后重跑）
% 第1章 周期结构

% 第1章 PERIODIC STRUCTURES

\chapter{周期结构}

\epigraphbook{再来一遍！再来一遍！再来一遍！}{Thomas Campbell}

\section{平移对称性}

如果大多数固体的最稳定结构不是规则的晶格，那么关于固体物理性质的理论实际上将无从建立。平移对称性（translational symmetry）的存在，使 $N$ 体问题被约化到可以处理的规模。这就是说，\emph{存在着基矢（basic vectors）$\mathbf{a}_1$、$\mathbf{a}_2$、$\mathbf{a}_3$，使得原子结构在经过任何一个矢量——它为这些基矢的整数倍之和——的平移后保持不变}。

\emph{实际上}，这只是一个理想。实验室样品必然尺寸有限，因此我们不能把结构一直延伸越过边界。但受此影响的区域只有表面附近的几层原子；在一块含 $N$ 个原子的样品中，它们只占大约 $N^{2/3}$ 个原子——就是说，在宏观样品中不过 $10^{8}$ 个原子里有一个。大多数晶体固体在结构上也是不完美的，存在缺陷、杂质和位错，扰乱着原子排列的规则性。这类不完整性引起许多有趣的物理现象，但在本书的讨论中，除非偶然涉及，我们将一概置之不理。我们在这里主要关注的，是完美的理想固体及其所表现的性质；而那些把固体仅仅当作些许污垢、微细裂纹和结构瑕疵的基体、载体或背景的现象，则属于另一个论域。

我们用空间点阵（space lattice）或 Bravais 网（Bravais net）来表示平移群。从某一点出发，构造出由基本平移可从它到达的所有点。这些点便是格点（lattice sites），由集合
\begin{equation*}
\boldsymbol{l} = l_1\mathbf{a}_1 + l_2\mathbf{a}_2 + l_3\mathbf{a}_3,
\tag{1.1}
\end{equation*}
给出，其中 $l_1$、$l_2$、$l_3$ 为整数（图 1）。

但固体是一种物理结构——并不是一组数学点。设想在我们所选定的原点 $O$ 附近有若干原子等等。平移不变性坚持要求：在每个格点周围，都必须恰好有完全相同的原子，以完全相同的方式配置（图 2）。

显然，只要指明单个晶胞（unit cell）的内容——例如由基矢 $\mathbf{a}_1$、$\mathbf{a}_2$、$\mathbf{a}_3$ 张成的平行六面体——我们就能确定整个晶体的物理布局。整个晶体就是这种物体的无穷多次重复，像砌墙的砖一样堆叠起来。不过，晶胞的具体取法在某种程度上是任意的：任何大小、形状与取向合适的平行六面体显然都行——正如我们在
\begin{figure}[H]
\centering
\includegraphics[width=0.72\textwidth]{fig_1.pdf}
\caption*{图 1\quad （原书此图无图注。）}
\end{figure}
\begin{figure}[H]
\centering
\includegraphics[width=0.72\textwidth]{fig_2.pdf}
\caption*{图 2\quad 可供选择的晶胞。}
\end{figure}
图 2 中看到的那样。或许不那么显而易见的是，晶胞的形状还可以在一定程度上加以改变。例如，假设结构中某一点具有中心对称性（因而在所有等价点处也都是如此）。这一点便很适合选作晶胞的中心，晶胞自身也具有中心对称性。系统地做到这一点的办法，是构造一个 Wigner–Seitz 原胞（Wigner–Seitz cell），也就是从选定的中心向最近邻的等价格点，画出各平移矢量的垂直平分面。全部平分面之内的体积显然是一个晶胞——它就是其中各点离选定中心比离任何其他格点都近的那个区域。

晶胞可以含有一个或多个原子。自然，如果只含一个原子，我们就把它置于格点中心，并说我们有一个 Bravais 格子（Bravais lattice）。如果每个晶胞含几个原子，那么我们就有了一个带基元的格子（lattice with a basis）。在后面的大部分内容中，我们将不加特别声明地假定结构是 Bravais 格子。这只是为简单起见；实际上只有少数几种单质固体——例如碱金属——才具有这种结构。

晶体学（crystallography）这门科学，关注的是对所有可能类型的晶体结构进行清点和分类，并确定实际晶体的实际结构。

\begin{figure}[H]
\centering
\includegraphics[width=0.72\textwidth]{fig_3.pdf}
\caption*{图 3\quad Wigner–Seitz 原胞。}
\end{figure}

结构按其对称性质分类，例如绕轴旋转的不变性、对平面反射的不变性等等。这些对称性常常对简化理论计算十分重要，并且在讨论定义固体宏观性质所需要的参数数目时，可以发挥很大的威力。然而，要充分利用这套理论，需要群论（group theory）这门数学，它会使我们偏离正题太远。如果我们主要局限于非常简单的固体，那么大多数对称性质无须正式的代数分析，单凭观察就能发现。无论如何，关于晶体学、以及关于群论及其在固体理论中的应用，都有许多优秀的著作。

在这些书中，各类 Bravais 格子等等都有详细的铺陈。我们在这里只考虑一个例子，它体现了本学科的许多原理，而且作为一种某些元素实际上采取的结构，也极为重要。这就是体心立方（body-centred cubic，B.C.C.）结构，如图 4 所示。

\begin{figure}[H]
\centering
\includegraphics[width=0.6\textwidth]{fig_4.pdf}
\caption*{图 4\quad 体心立方格子。(a)立方晶胞。(b)Bravais 格子的生成矢。(c)Wigner–Seitz 原胞。}
\end{figure}

\begin{figure}[H]
\centering
\includegraphics[width=0.72\textwidth]{fig_5.pdf}
\caption*{图 5\quad B.C.C. 格子 Wigner–Seitz 原胞的堆垛。}
\end{figure}

乍一看，这是一个每个立方晶胞含两个原子的格子，或者说是两个相互穿插的简立方亚格子，由
\begin{equation*}
\left.
\begin{aligned}
\boldsymbol{l} &= l_x\mathbf{a}_x + l_y\mathbf{a}_y + l_z\mathbf{a}_z, \\
\boldsymbol{l}' &= \left(l_x + \tfrac{1}{2}\right)\mathbf{a}_x + \left(l_y + \tfrac{1}{2}\right)\mathbf{a}_y + \left(l_z + \tfrac{1}{2}\right)\mathbf{a}_z,
\end{aligned}
\right\}
\tag{1.2}
\end{equation*}
定义，其中 $l_x$、$l_y$、$l_z$ 均为整数。但如果写下
\begin{equation*}
\left.
\begin{aligned}
\mathbf{a}_1 &= \tfrac{1}{2}\left(-\mathbf{a}_x + \mathbf{a}_y + \mathbf{a}_z\right), \\
\mathbf{a}_2 &= \tfrac{1}{2}\left(\mathbf{a}_x - \mathbf{a}_y + \mathbf{a}_z\right), \\
\mathbf{a}_3 &= \tfrac{1}{2}\left(\mathbf{a}_x + \mathbf{a}_y - \mathbf{a}_z\right),
\end{aligned}
\right\}
\tag{1.3}
\end{equation*}
我们便能由
\begin{equation*}
\boldsymbol{l} = l_1\mathbf{a}_1 + l_2\mathbf{a}_2 + l_3\mathbf{a}_3
\tag{1.4}
\end{equation*}
（$l_1$、$l_2$、$l_3$ 为整数）生成两个亚格子的全部格点。我们发现，落在立方体的中心还是顶角上，视 $(l_1+l_2+l_3)$ 为奇数还是偶数而定。

因此（图 4(b)），这其实是一个 Bravais 格子。我们可以不用立方晶胞，而改用 Wigner–Seitz 原胞（图 4(c)）；它的构造方法是，沿从中心到角点的对角线在半程处把立方体的各个角截去。这个图形显然具有与立方体相同的对称性——例如，原来的矢量 $\mathbf{a}_x$、$\mathbf{a}_y$、$\mathbf{a}_z$ 都是四重对称轴。它还（或许比原来的立方体更清楚地）表明，矢量 $\mathbf{a}_1$、$\mathbf{a}_2$、$\mathbf{a}_3$，即立方体的对角线，是三重旋转对称轴，因为它们穿过原胞的六边形面。设想一下这些多面体如何能堆砌成原来的格子（图 5），是一次很好的练习。

\section{周期函数}

要为一个晶体结构定义物理模型，我们需要给空间中的某个函数 $f(\mathbf{r})$ 赋值——局域电子密度、静电势等等，它们应能被辨认为原子的一种排列（例如图 6）。我们关于平移对称性的断言意味着，这必须是一个\emph{多重周期函数}（multiply periodic function）
\begin{equation*}
f(\mathbf{r}+\boldsymbol{l}) = f(\mathbf{r})
\tag{1.5}
\end{equation*}
对空间中所有的点 $\mathbf{r}$、对所有的格平移（lattice translations）都成立。

\begin{figure}[H]
\centering
\includegraphics[width=0.72\textwidth]{fig_6.pdf}
\caption*{图 6\quad 多重周期函数。}
\end{figure}

\begin{figure}[H]
\centering
\includegraphics[width=0.72\textwidth]{fig_7.pdf}
\caption*{图 7\quad （原书此图无图注。）}
\end{figure}

在一维情形，我们对周期函数再熟悉不过了。例如，如图 7 那样，可以有
\begin{equation*}
f(x+l) = f(x),
\tag{1.6}
\end{equation*}
其中 $l$ 取 $l_1 a$ 的形式，$l_1$ 为整数，而 $a$ 是函数的周期。我们还知道，任何这类函数都可以表示成 Fourier 级数，
\begin{equation*}
f(x) = \sum_n A_n e^{2\pi i n x/a},
\tag{1.7}
\end{equation*}
其中 $n$ 是整数。让我们把它写成
\begin{equation*}
f(x) = \sum_g A_g e^{igx},
\tag{1.8}
\end{equation*}
其中的量 $g$ 属于倒格子长度（reciprocal lattice lengths）\footnote{在这个定义中把因子 $2\pi$ 包括进来是方便的；虽然在晶体学中惯例是把 $2\pi$ 明确地留在指数上。}的集合
\begin{equation*}
g_n = n\,\frac{2\pi}{a}.
\tag{1.9}
\end{equation*}

大家熟知，(1.8) 中的系数由
\begin{equation*}
A_g = \frac{1}{a}\int_{\text{cell}} f(x)\, e^{-igx}\, dx,
\tag{1.10}
\end{equation*}
定义；在这个例子里，积分范围譬如取 $0 < x \leqslant a$——它只需是晶格的一个晶胞。(1.8) 蕴含 (1.6) 这一初等证明，可以由如下条件导出：对任何 $g$ 值，
\begin{equation*}
e^{igl} = 1
\tag{1.11}
\end{equation*}
对所有平移 $l$ 成立。当 $g$ 取其一个允许值 $g_n$ 而 $l = l_1 a$ 时，这不过是说
\begin{equation*}
gl = n\,\frac{2\pi}{a}\, l_1 a = n l_1\, 2\pi = \text{integer} \times 2\pi
\tag{1.12}
\end{equation*}
罢了。(1.10) 的推导同样可以由这同一个关系得出。我们在这里不关心数学上病态的函数，尽可以放心地随意使用朴素的 Fourier 理论。

把这个定理推广到具有三条正交轴的结构相当容易。设这些轴为 $\mathbf{a}_x$、$\mathbf{a}_y$、$\mathbf{a}_z$，那么，只要
\begin{equation*}
f(\mathbf{r}+\boldsymbol{l}) \equiv f(\mathbf{r}+l_x\mathbf{a}_x+l_y\mathbf{a}_y+l_z\mathbf{a}_z) = f(\mathbf{r}),
\tag{1.13}
\end{equation*}
便有
\begin{equation*}
f(\mathbf{r}) = \sum_{g_x,g_y,g_z} A(g_x,g_y,g_z) \exp\left\{ i(g_x x + g_y y + g_z z)\right\},
\tag{1.14}
\end{equation*}
其中 $g_x$ 等等各自都是从集合 $2\pi n/a_x$ 等等中取的一个倒格子长度。证明可以这样做：先把 $f(\mathbf{r})$ 沿 $x$ 方向展开成 Fourier 级数，然后证明该级数的每个系数都是 $y$ 的周期函数，因而可以再展开成一个级数——对 $z$ 亦照此办理。

让我们把 (1.14) 改写成
\begin{equation*}
f(\mathbf{r}) = \sum_{\boldsymbol{g}} A_{\boldsymbol{g}} e^{i\boldsymbol{g}\cdot\mathbf{r}},
\tag{1.15}
\end{equation*}
其中 $\boldsymbol{g}$ 是分量为 $(g_x,g_y,g_z)$ 的矢量。这种矢量具有 (1.11) 所表达的性质，也就是说，对任何一个这样的矢量
\begin{align*}
\boldsymbol{g}\cdot\boldsymbol{l} &= \left(g_x l_x \mathbf{a}_x + g_y l_y \mathbf{a}_y + g_z l_z \mathbf{a}_z\right) \\
&= \frac{2\pi n_x}{a_x} l_x \mathbf{a}_x + \frac{2\pi n_y}{a_y} l_y \mathbf{a}_y + \frac{2\pi n_z}{a_z} l_z \mathbf{a}_z \\
&= 2\pi \times \text{integer},
\tag{1.16}
\end{align*}
无论 $\boldsymbol{l}$ 取什么值。于是，
\begin{equation*}
e^{i\boldsymbol{g}\cdot\boldsymbol{l}} = 1
\tag{1.17}
\end{equation*}
对一切格矢 $\boldsymbol{l}$、对一切倒格矢（reciprocal lattice vector）$\boldsymbol{g}$ 都成立。

显然，这个条件足以使级数 (1.15) 表示一个像 (1.5) 那样具有晶格周期性的函数：
\begin{align*}
f(\mathbf{r}+\boldsymbol{l}) &= \sum_{\boldsymbol{g}} A_{\boldsymbol{g}}\, e^{i\boldsymbol{g}\cdot(\mathbf{r}+\boldsymbol{l})} = \sum_{\boldsymbol{g}} A_{\boldsymbol{g}}\, e^{i\boldsymbol{g}\cdot\mathbf{r}}\, e^{i\boldsymbol{g}\cdot\boldsymbol{l}} \\
&= \sum_{\boldsymbol{g}} A_{\boldsymbol{g}}\, e^{i\boldsymbol{g}\cdot\mathbf{r}} = f(\mathbf{r}).
\tag{1.18}
\end{align*}
人们不难设计出一个证明，说明它也是必要条件，也就是说，级数 (1.15) 中只能含有 $g$ 值满足条件 (1.17) 的那些项。

剩下的只是为非正交格子构造倒格矢。这很容易做到：取基矢的倒格子三矢组（reciprocal triad），即
\begin{equation*}
\mathbf{b}_1 = \frac{\mathbf{a}_2\wedge\mathbf{a}_3}{\mathbf{a}_1\cdot\mathbf{a}_2\wedge\mathbf{a}_3}, \qquad
\mathbf{b}_2 = \frac{\mathbf{a}_3\wedge\mathbf{a}_1}{\mathbf{a}_1\cdot\mathbf{a}_2\wedge\mathbf{a}_3}, \qquad
\mathbf{b}_3 = \frac{\mathbf{a}_1\wedge\mathbf{a}_2}{\mathbf{a}_1\cdot\mathbf{a}_2\wedge\mathbf{a}_3},
\tag{1.19}
\end{equation*}
并写出
\begin{align*}
\boldsymbol{g} &= g_1\mathbf{b}_1 + g_2\mathbf{b}_2 + g_3\mathbf{b}_3 \\
&= 2\pi n_1\mathbf{b}_1 + 2\pi n_2\mathbf{b}_2 + 2\pi n_3\mathbf{b}_3,
\tag{1.20}
\end{align*}
其中 $n_1$ 等等为整数。

由初等矢量分析即可验证：$\mathbf{b}_1\cdot\mathbf{a}_1 = 1$ 等等，以及 $\mathbf{b}_1\cdot\mathbf{a}_2 = 0$ 等等；于是
\begin{align*}
\boldsymbol{g}\cdot\boldsymbol{l} &= \left(g_1\mathbf{b}_1 + g_2\mathbf{b}_2 + g_3\mathbf{b}_3\right)\cdot\left(l_1\mathbf{a}_1 + l_2\mathbf{a}_2 + l_3\mathbf{a}_3\right) \\
&= g_1 l_1 + g_2 l_2 + g_3 l_3 \\
&= 2\pi \times \text{integer}.
\tag{1.21}
\end{align*}
因此，集合 (1.20) 中的每一个矢量都满足条件 (1.17)。

我们还可以把公式 (1.10) 加以推广，用来求级数 (1.15) 中的每个系数。它必须是一个体积分，我们把它写作
\begin{equation*}
A_{\boldsymbol{g}} = \frac{1}{v_{\text{cell}}} \int_{\text{cell}} f(\mathbf{r})\, e^{-i\boldsymbol{g}\cdot\mathbf{r}}\,\mathrm{d}\mathbf{r}.
\tag{1.22}
\end{equation*}

%（以下为 §1.2 的收尾段落，印于书页 9 顶部，接式 (1.22) 之后）
证明只需把整个级数乘以 $\exp(-i\boldsymbol{g}\cdot\mathbf{r})$，然后积分。若 $\boldsymbol{g}$ 具有 (1.20) 的形式，则 $\exp(i\boldsymbol{g}\cdot\mathbf{r})$ 在一个晶胞上的积分显然为零，除非 $\boldsymbol{g}=0$。

\section{倒格子的性质}

由 (1.20) 定义的矢量，即
\begin{equation*}
\boldsymbol{g} = n_1\cdot 2\pi\mathbf{b}_1 + n_2\cdot 2\pi\mathbf{b}_2 + n_3\cdot 2\pi\mathbf{b}_3,
\tag{1.23}
\end{equation*}
生成一个点阵，其晶胞由矢量 $2\pi\mathbf{b}_1$、$2\pi\mathbf{b}_2$、$2\pi\mathbf{b}_3$ 张成。这就是我们原来正格子（direct lattice）的倒格子（reciprocal lattice）。它是一个不变的几何客体，其性质是固体理论的基础。它的几个初等几何性质很容易推导。

(i) \emph{倒格子的每一个矢量都垂直于正格子的一族点阵面（lattice planes）。}

任取一个倒格矢 $\boldsymbol{g}$ 和一个格矢 $\boldsymbol{l}$，考虑关系式 (1.21)：
\begin{align*}
\boldsymbol{g}\cdot\boldsymbol{l} &= 2\pi(n_1 l_1 + n_2 l_2 + n_3 l_3) \\
&= 2\pi N,
\tag{1.24}
\end{align*}
其中 $N$ 是整数。这告诉我们，矢量 $\boldsymbol{l}$ 在 $\boldsymbol{g}$ 方向上的投影长度为
\begin{equation*}
d = \frac{2\pi N}{|\boldsymbol{g}|}.
\tag{1.25}
\end{equation*}
但是，具有这一性质的格点有无穷多个。例如，设 $\boldsymbol{l}'$ 是由整数
\begin{equation*}
l'_1 = l_1 - mn_3, \qquad l'_2 = l_2 - mn_3, \qquad l'_3 = l_3 + m(n_1+n_2),
\tag{1.26}
\end{equation*}
所代表的格点，其中 $m$ 是整数（译注：原书第三式误印作 $l_1+m(n_1+n_2)$，按上下文应为 $l_3$）。显然
\begin{equation*}
\boldsymbol{g}\cdot\boldsymbol{l}' = \boldsymbol{g}\cdot\boldsymbol{l} = 2\pi N.
\tag{1.27}
\end{equation*}
于是 $\boldsymbol{l}'$ 在 $\boldsymbol{g}$ 上的投影与 $\boldsymbol{l}$ 相同，因而必定也位于与 $\boldsymbol{g}$ 垂直、距原点为 $d$ 的平面上。这样，若这个平面上有一个格点，其上就有无穷多个格点；我们便构造出了一个\emph{点阵面}。正格子与倒格子之间的这种关系，当然是三维几何中人所共知的点与平面对偶性的一个特例。

(ii) \emph{若 $\boldsymbol{g}$ 的分量没有公因子，则 $|\boldsymbol{g}|$ 与垂直于 $\boldsymbol{g}$ 的点阵面的间距成反比。}

这由式 (1.25) 可得（见图 8）。若 $(n_1, n_2, n_3)$ 没有公因子，则总能找到一个格矢 $\boldsymbol{l}''$，其分量满足
\begin{equation*}
\boldsymbol{g}\cdot\boldsymbol{l}'' = 2\pi(N+1).
\tag{1.28}
\end{equation*}
这意味着，包含 $\boldsymbol{l}''$ 的点阵面位于距原点
\begin{equation*}
d'' = \frac{2\pi(N+1)}{|\boldsymbol{g}|}
\tag{1.29}
\end{equation*}
处——亦即它与包含 $\boldsymbol{l}$ 的点阵面相距 $2\pi/|\boldsymbol{g}|$。

\begin{figure}[H]
\centering
\includegraphics[width=0.72\textwidth]{fig_8.pdf}
\caption*{图 8\quad （原书此图无图注。）}
\end{figure}

从这两个初等几何结果可以看出，表征点阵诸平面最简单的方式就是用它们的法线，把它们表示为倒格矢。正格子中最显著的平面，是格点最密集的那些平面。由于格点在空间中的密度是常数，最显著的平面必定是相隔最远的那些平面，即具有最小倒格矢的那些平面。

用相应的倒格矢来标示点阵面，等价于经典结晶学中所用的米勒指数（Miller indices）。设有一个点阵面，法线为 $\boldsymbol{g}$，其上所有的点 $\boldsymbol{l}$ 都满足 (1.24)。取 $l_2 = l_3 = 0$ 的点，则得 $l_1 = N/n_1$，于是这个平面沿 $\mathbf{a}_1$ 轴的截距长度为
\begin{equation*}
d_1 = \left(\frac{N}{n_1}\right) a_1.
\tag{1.30}
\end{equation*}
类似地，这个平面与 $\mathbf{a}_2$ 轴相交于距原点
\begin{equation*}
d_2 = \left(\frac{N}{n_2}\right) a_2
\tag{1.31}
\end{equation*}
处。该平面沿各轴的截距，以相应基矢的长度为单位来量度，与整数 $n_1$、$n_2$、$n_3$ 成反比。除去公因子之后，这些整数就是这个点阵面的米勒指数，写成 $(n_1, n_2, n_3)$ 的形式。

显然，格点密集的面——即米勒指数小的面——最容易在天然晶体中出现，无论是在生长时还是解理之后。研究这类晶面的几何曾是经典结晶学的精髓，而其
\begin{figure}[H]
\centering
\includegraphics[width=0.72\textwidth]{fig_9.pdf}
\caption*{图 9\quad 具有大米勒指数的晶面，其间距比主对称面的间距更小。}
\end{figure}
关键性的发现是：可以找到单独一组矢量 $\mathbf{a}_1$、$\mathbf{a}_2$、$\mathbf{a}_3$，使得宏观晶体上所有观察到的晶面都能用小的米勒指数来表示。

应当提到两种记号约定。符号 $(1\bar{1}\bar{1})$ 代表 $(1,-1,-1)$ 这组平面；为简洁起见，把负号写在数字上方。带花括号的符号 $\{110\}$ 代表在对称性上与 $(110)$ 等价的各组不同的平面。因此，对立方结构而言，它可以包括 $(101)$、$(011)$、$(\bar{1}\bar{1}0)$、$(1\bar{1}0)$，等等。

(iii) \emph{倒格子晶胞的体积与正格子晶胞的体积成反比。}

这用初等矢量分析即可得出。倒格子晶胞由 $2\pi\mathbf{b}_1$、$2\pi\mathbf{b}_2$、$2\pi\mathbf{b}_3$ 张成。利用 (1.19)，其体积为
\begin{align*}
(2\pi)^3(\mathbf{b}_1\cdot\mathbf{b}_2\wedge\mathbf{b}_3)
&= (2\pi)^3
\frac{(\mathbf{a}_2\wedge\mathbf{a}_3)\cdot\{(\mathbf{a}_3\wedge\mathbf{a}_1)\wedge(\mathbf{a}_1\wedge\mathbf{a}_2)\}}
{(\mathbf{a}_1\cdot\mathbf{a}_2\wedge\mathbf{a}_3)^3} \\
&= (2\pi)^3
\frac{(\mathbf{a}_2\wedge\mathbf{a}_3)\cdot[\{\mathbf{a}_3\cdot(\mathbf{a}_1\wedge\mathbf{a}_2)\}\mathbf{a}_1 - \{\mathbf{a}_1\cdot(\mathbf{a}_1\wedge\mathbf{a}_2)\}\mathbf{a}_3]}
{(\mathbf{a}_1\cdot\mathbf{a}_2\wedge\mathbf{a}_3)^3} \\
&= \frac{(2\pi)^3}{(\mathbf{a}_1\cdot\mathbf{a}_2\wedge\mathbf{a}_3)} \\
&= \frac{8\pi^3}{v_c},
\tag{1.32}
\end{align*}
其中 $v_c$ 是由 $\mathbf{a}_1$、$\mathbf{a}_2$、$\mathbf{a}_3$ 张成的晶胞的体积。

这里出现因子 $8\pi^3$，是由于我们对倒格子的定义方式。更常见的约定是写成
\begin{equation*}
e^{i\boldsymbol{g}\cdot\boldsymbol{l}} \equiv \exp(2\pi i\,\mathbf{K}_g\cdot\mathbf{R}_l),
\tag{1.33}
\end{equation*}
其中 $\mathbf{R}_l$ 是格矢，而矢量 $\mathbf{K}_g$ 定义为晶胞体积为 $1/v_c$ 的另一个倒格子的矢量。这种记法的缺点是因子 $2\pi$ 总是在指数中冒出来；它也同自由电子波函数的惯用符号
\begin{equation*}
\psi = e^{i\mathbf{k}\cdot\mathbf{r}}
\tag{1.34}
\end{equation*}
不相协调。

(iv) \emph{正格子正是它自己的倒格子的倒格子。}

这已隐含在名称之中，也可以通过构造例如矢量 $(\mathbf{b}_1\wedge\mathbf{b}_2)/(\mathbf{b}_1\cdot\mathbf{b}_2\wedge\mathbf{b}_3)$ 并证明它与 $\mathbf{a}_3$ 全同来验证。考察 (1.32) 的推导过程也能看出这一点。

(v) \emph{倒格子的晶胞不必是平行六面体。}事实上，我们几乎总是在同倒格子的 Wigner--Seitz 胞（Wigner--Seitz cell）打交道。它叫做布里渊区（Brillouin zone）。

作为倒格子诸性质的一个例子，考虑一种重要且常见的结构——面心立方（face-centred cubic）格子（f.c.c. 结构）。它由四个相互穿插的简单立方格子搭成；随便看其中哪一个，都会看到它的晶胞每个面的中心各有一个原子，同时立方体的角上也分布着格点（图 10(a)）。

乍一看，这像一个带基元的格子：由 $(\mathbf{a}_x,\mathbf{a}_y,\mathbf{a}_z)$ 生成的立方晶胞里有四个原子。

但是，如果改取立方体各个面的半对角线
\begin{equation*}
\left.
\begin{aligned}
\mathbf{a}_1 &= (0, \tfrac{1}{2}a, \tfrac{1}{2}a), \qquad
\mathbf{a}_2 = (\tfrac{1}{2}a, 0, \tfrac{1}{2}a), \\
\mathbf{a}_3 &= (\tfrac{1}{2}a, \tfrac{1}{2}a, 0),
\end{aligned}
\right\}
\tag{1.35}
\end{equation*}
\begin{figure}[H]
\centering
\includegraphics[width=0.58\textwidth]{fig_10.pdf}
\caption*{图 10\quad 面心立方格子。(a) 四个相互穿插的子晶格。(b) 作为 Bravais 格子。}
\end{figure}

那么用这些矢量的倍数的组合就能到达任何一个格点（图 10(b)）。因此，f.c.c. 格子确实是一个 Bravais 格子，$\mathbf{a}_1$、$\mathbf{a}_2$、$\mathbf{a}_3$ 就是它的生成矢。

现在我们本可以用代数方法构造出倒格子来，但利用上面指出的几何性质更为容易。例如，显然存在垂直于立方体棱边的重要点阵面。相对沿 $\mathbf{a}_x$、$\mathbf{a}_y$、$\mathbf{a}_z$ 方向的通常笛卡儿轴而言，这些面的米勒指数必为 (100)、(010)、(001) 等，即它们构成 $\{100\}$ 组。还可以看出，这些面彼此相距 $\tfrac{1}{2}a$。于是，它们对应于长度为
\begin{equation*}
|\boldsymbol{g}| = 2\pi/\tfrac{1}{2}a
\end{equation*}
沿倒格子空间中直角笛卡儿轴方向的倒格矢。我们的倒格子必须包含由这些矢量生成的整个简单立方格子。

还有一种重要的点阵面垂直于立方晶胞的对角线。这种面在三条轴上的截距相等，故其指数必属 $\{111\}$ 组。这些面
\begin{figure}[H]
\centering
\includegraphics[width=0.72\textwidth]{fig_11.pdf}
\caption*{图 11\quad (a) f.c.c. 格子的 {111} 面。(b) 相应的倒格子晶胞。}
\end{figure}
彼此相隔的距离等于立方体整条对角线的三分之一，即 $a/\sqrt{3}$。因此，相应的倒格矢有三个相等的分量，总长为 $2\pi\sqrt{3}/a$；它必是
\begin{equation*}
\boldsymbol{g}' = \frac{2\pi}{a}(1, 1, 1).
\tag{1.36}
\end{equation*}
在倒格子空间中，它显然沿立方晶胞对角线的方向——但长度只有对角线的一半。于是，它生成简单立方格子的各个体心。

其余的点阵面现在都可以如法研究——例如 $\{110\}$ 组面，它们垂直于基本立方体各面的对角线——不过与它们相应的倒格矢全都已经包含在由 $\{100\}$ 面和 $\{111\}$ 面生成的那些点之中。因此，f.c.c. 格子的倒格子是 b.c.c.（当然，反之亦然）。

f.c.c.格子的布里渊区就是这个倒格子的Wigner–Seitz原胞。这个图形我们已经在图 4(c)中见过——一个兼有正方形面与六边形面的多面体。正方形面的中心对应于正格子诸立方晶面法线的方向；六边形面的中心则是诸对角面法线的方向。

还有一点颇有意思：f.c.c.结构是一种“密堆积”（close-packed）结构。依次叠合一层层(111)面就可以把它构造出来，每一层都是一个六边形网络，就像由刚性球排成的那样。这也是金属常常采取这种结构的原因之一：金属的内聚并不依赖于强方向性的键，而在很大程度上是一种体积效应。

\section{Bloch 定理}

我们已经学会了如何表示具有晶格周期性的函数。但对于一种物理理论来说这些还不够；我们还需要考虑结构的各种激发，它们会破坏严格的平移对称性。这类激发有若干种，其中最重要的是格波（lattice waves），即原子绕其平衡位置的振动；以及电子态（electron states），即电子在静态晶格的场中的运动。在第10章我们还将考虑自旋波（spin waves），即定域在晶体各原子上的那些自旋的激发。

所有这些激发都以一个动力学方程——或一个Schrödinger方程，或一个自旋交换哈密顿量——为特征，它们\emph{在晶格平移下不变}。于是，若 $\mathcal{V}(\mathbf{r})$ 是位于 $\mathbf{r}$ 处的电子所感受到的势，则对一切 $\mathbf{l}$ 都有 $\mathcal{V}(\mathbf{r}+\mathbf{l})=\mathcal{V}(\mathbf{r})$。电子的波函数 $\psi(\mathbf{r})$ 必须满足Schrödinger方程
\begin{equation*}
\left(-\frac{\hbar^{2}}{2m}\nabla^{2}+\mathcal{V}(\mathbf{r})-\mathcal{E}\right)\psi=0,
\tag{1.37}
\end{equation*}
在把作用于 $\psi$ 的算符中的 $\mathbf{r}$ 换成 $\mathbf{r}+\mathbf{l}$ 之后，这个方程保持不变。

在格波和自旋波的情形，形式体系要更复杂一些，但原理相同。设
\begin{equation*}
(u_1,\,u_2,\,\dots,\,u_{\mathbf{n}},\,\dots)
\end{equation*}
是位于 $1,2,\dots,\mathbf{n},\dots$ 各格点上的原子的位移（或自旋位移，或自旋偏离算符，或随便什么东西）。运动方程（或哈密顿量）将把这些变量（或算符）包含在内，其方式是：如果我们作一个晶格平移 $\mathbf{l}$——也就是说，把公式里原先出现 $u_{\mathbf{n}}$ 的地方统统换成 $u_{\mathbf{n}+\mathbf{l}}$——我们将看不出任何差别。换言之，晶格的所有晶胞都等价而且不可分辨。这颇像所谓“宇宙学原理”（cosmological principle）：无论我们从哪一点观察，宇宙看起来本质上都是一样的。

为了把这个性质形式化，我们采用如下记号。令 $\mathcal{H}(\mathbf{0})$ 和 $|0\rangle$ 表示晶格平移之前的哈密顿算符和波函数；令 $\mathcal{H}(\mathbf{l})$ 和 $|\mathbf{l}\rangle$ 表示平移 $\mathbf{l}$ 之后的同样一些数学对象。于是，对电子而言 $|0\rangle\equiv\psi(\mathbf{r})$，$|\mathbf{l}\rangle\equiv\psi(\mathbf{r}+\mathbf{l})$。对格波而言，若 $|0\rangle$ 表示位于 $\mathbf{n}$ 处的原子正经历某一特定运动的态，则 $|\mathbf{l}\rangle$ 就表示位于 $(\mathbf{n}+\mathbf{l})$ 处的原子正作同一运动的态。

平移不变性的表述便是
\begin{equation*}
\mathcal{H}(\mathbf{l})=\mathcal{H}(\mathbf{0}).
\tag{1.38}
\end{equation*}
现在，本征态必须满足如下形式的方程\footnote{在格波的经典运动方程的情形，频率 $\omega$ 的平方扮演能量的角色，而“哈密顿量”则是在我们假定所处理的是对平衡的小偏离时导出的那组联立线性微分方程的矩阵。参看 §2.1。}
\begin{equation*}
\mathcal{H}(\mathbf{0})\,|0\rangle=\mathcal{E}\,|0\rangle.
\tag{1.39}
\end{equation*}
这个方程与
\begin{equation*}
\mathcal{H}(\mathbf{l})\,|\mathbf{l}\rangle=\mathcal{E}\,|\mathbf{l}\rangle,
\tag{1.40}
\end{equation*}
完全相同，后者只不过是把算符与态函数中的所有变量作一次形式上的重新标记而已。但由式 (1.38)，这意味着
\begin{equation*}
\mathcal{H}(\mathbf{0})\,|\mathbf{l}\rangle=\mathcal{E}\,|\mathbf{l}\rangle,
\tag{1.41}
\end{equation*}
也就是说，态函数 $|\mathbf{l}\rangle$ 也是 $|0\rangle$ 所满足的那个方程的解。由于 $|\mathbf{l}\rangle$ 未必与 $|0\rangle$ 恒同，看上去我们似乎只要把碰巧找到的任何一个解加以平移，就能生成运动方程的数不清的解。当然，所有这些解在能量上都是简并的。

这显然是荒谬的；新解必定都以某种方式与我们原来的解等价。有两种情形需要考虑：

(i) 设 $|0\rangle$ 确实是\emph{非简并}的（这种情况实际上绝不会发生，因为一切晶格除平移不变性之外还具有镜面对称性）。那么唯一的可能是：每个新态 $|\mathbf{l}\rangle$ 都是 $|0\rangle$ 的某个倍数，从而在物理上与它不可分辨。必定存在某个数 $\lambda_1$，使得
\begin{equation*}
|\mathbf{a}_1\rangle=\lambda_1|0\rangle
\tag{1.42}
\end{equation*}
这是沿 $\mathbf{a}_1$ 方向走一步的结果。归一化要求
\begin{equation*}
|\lambda_1|^{2}=1,
\tag{1.43}
\end{equation*}
于是可以取
\begin{equation*}
\lambda_1=e^{ik_1},
\tag{1.44}
\end{equation*}
其中 $k_1$ 为实数。类似地，对其余两个基本方向上的单位平移，可以取
\begin{equation*}
|\mathbf{a}_2\rangle=e^{ik_2}|0\rangle;\qquad|\mathbf{a}_3\rangle=e^{ik_3}|0\rangle,
\tag{1.45}
\end{equation*}
于是一般平移必须导致如下关系
\begin{align*}
|\mathbf{l}\rangle&\equiv|l_1\mathbf{a}_1+l_2\mathbf{a}_2+l_3\mathbf{a}_3\rangle\notag\\
&=e^{ik_1}|(l_1-1)\mathbf{a}_1+l_2\mathbf{a}_2+l_3\mathbf{a}_3\rangle\notag\\
&=e^{ik_1 l_1}|0+l_2\mathbf{a}_2+l_3\mathbf{a}_3\rangle\notag\\
&=e^{i(k_1 l_1+k_2 l_2+k_3 l_3)}|0\rangle
\tag{1.46}
\end{align*}
这里我们沿 $\mathbf{a}_1$ 方向走 $l_1$ 步、沿 $\mathbf{a}_2$ 方向走 $l_2$ 步，如此等等，每走一步都乘上相应的因子。

我们按如下形式定义一个矢量 $\mathbf{k}$：
\begin{equation*}
\mathbf{k}=k_1\mathbf{b}_1+k_2\mathbf{b}_2+k_3\mathbf{b}_3,
\tag{1.47}
\end{equation*}
其中 $\mathbf{b}_1,\mathbf{b}_2,\mathbf{b}_3$ 是 $\mathbf{a}_1,\mathbf{a}_2,\mathbf{a}_3$ 的倒三基矢（reciprocal triad），如式 (1.19) 中那样。关系式 (1.46) 便成为
\begin{equation*}
|\mathbf{l}\rangle=e^{i\mathbf{k}\cdot\mathbf{l}}|0\rangle,
\tag{1.48}
\end{equation*}
这正是我们想要证明的结果：

\emph{对于任何一个满足Schrödinger方程（或其经典或量子的对应形式）的波函数/态函数，都存在一个矢量 $\mathbf{k}$，使得平移一个格矢 $\mathbf{l}$ 等价于乘上相位因子 $\exp(i\mathbf{k}\cdot\mathbf{l})$。}

每个不同的态函数可以有不同的波矢（wave-vector）$\mathbf{k}$，但每一个都必须满足这个条件。晶格具有平移不变性这个强的初始前提，对元激发的形式强加了一条强的限制条件。

但我们还需要在更一般的情形下给出证明。

(ii) 设 $|0\rangle$ 是简并的。为简单起见，设它是二重简并的，于是肯定存在两个能量相同而又彼此不同的态 $|0\rangle_1$ 和 $|0\rangle_2$。那么，晶格平移 $\mathbf{a}_1$ 至多能产生这两个态本身的线性组合。于是
\begin{equation*}
\left.\begin{aligned}
|\mathbf{a}_1\rangle_1&=\mathsf{T}_1^{11}\,|0\rangle_1+\mathsf{T}_1^{12}\,|0\rangle_2,\\
|\mathbf{a}_1\rangle_2&=\mathsf{T}_1^{21}\,|0\rangle_1+\mathsf{T}_1^{22}\,|0\rangle_2,
\end{aligned}\right\}
\tag{1.49}
\end{equation*}
诸数 $\mathsf{T}_1^{11}$ 等之所以这样书写，是因为它们是矩阵 $\mathsf{T}_1$ 的元素；为了保持归一化，$\mathsf{T}_1$ 必须是幺正矩阵（unitary matrix）。为简洁起见，把这对方程写成
\begin{equation*}
\left(\,|\mathbf{a}_1\rangle\right)=\mathsf{T}_1\left(\,|0\rangle\right).
\tag{1.50}
\end{equation*}
但 $|0\rangle_1$ 和 $|0\rangle_2$ 既然是简并的，就不是唯一定义的。我们原本可以从任何另外两个态出发，只要它们是这两个态的线性组合。例如，我们也许可以从
\begin{equation*}
\left.\begin{aligned}
|0\rangle_1&=\mathsf{S}^{11}\,|0\rangle_1+\mathsf{S}^{12}\,|0\rangle_2,\\
|0\rangle_2&=\mathsf{S}^{21}\,|0\rangle_1+\mathsf{S}^{22}\,|0\rangle_2,
\end{aligned}\right\}
\tag{1.51}
\end{equation*}
出发，其中诸数 $\mathsf{S}^{11}$ 等是另一个（任意的）幺正矩阵 $\mathsf{S}$ 的元素。

现在我们这样选取 $\mathsf{S}$，使矩阵 $\mathsf{S}\mathsf{T}_1\mathsf{S}^{-1}$ 成为对角矩阵。按标准的代数理论，这总是做得到的。例如，设它约化为如下形式
\begin{equation*}
\mathsf{S}\mathsf{T}_1\mathsf{S}^{-1}=
\begin{pmatrix}
e^{ik_1} & 0\\
0 & e^{ik_1'}
\end{pmatrix}.
\tag{1.52}
\end{equation*}
很容易验证，态 $|0\rangle_1$ 和 $|0\rangle_2$ 现在在平移之下的变换方式，恰好就像它们是非简并的一样，即如式 (1.42–1.45) 中那样：
\begin{equation*}
|\mathbf{a}_1\rangle_1=e^{ik_1}|0\rangle_1,\qquad|\mathbf{a}_1\rangle_2=e^{ik_1'}|0\rangle_2.
\tag{1.53}
\end{equation*}
态 $|0\rangle_1$ 有一个波矢，其在 $\mathbf{a}_1$ 方向上的分量为 $k_1$；态 $|0\rangle_2$ 对同一平移方向则有一个分量为 $k_1'$（一般说来与之不同）的波矢。

但是其他方向的平移又如何呢？考虑 $\mathbf{a}_2$。对初始态 $|0\rangle_1$ 和 $|0\rangle_2$，将有另一个矩阵 $\mathsf{T}_2$ 确定这一平移的结果，即像式 (1.49) 和 (1.50) 中那样符号式地写成
\begin{equation*}
\left(\,|\mathbf{a}_2\rangle\right)=\mathsf{T}_2\left(\,|0\rangle\right).
\tag{1.54}
\end{equation*}
现在我们也可以把这个矩阵 $\mathsf{T}_2$ 化成对角形式，正如式 (1.52) 中那样，只要选好一组适当的起始函数。乍一看，这似乎与我们用来使 $\mathsf{T}_1$ 对角化的那组函数不相容。但请考虑如下两种可供选择的约化：
\begin{equation*}
\left(\,|\mathbf{a}_1+\mathbf{a}_2\rangle\right)=
\left\{\begin{aligned}
\mathsf{T}_2(\,|\mathbf{a}_1\rangle)&=\mathsf{T}_2\mathsf{T}_1(\,|0\rangle),\\
\mathsf{T}_1(\,|\mathbf{a}_2\rangle)&=\mathsf{T}_1\mathsf{T}_2(\,|0\rangle).
\end{aligned}\right\}
\tag{1.55}
\end{equation*}
矩阵 $\mathsf{T}_1$ 和 $\mathsf{T}_2$ 由此彼此对易。矩阵代数中有一条定理告诉我们：这时存在一个幺正矩阵 $\mathsf{S}$，使得 $\mathsf{T}_1$ 和 $\mathsf{T}_2$ \emph{二者}都能\emph{同时}约化为对角形式。于是，态 $|0\rangle_1$ 和 $|0\rangle_2$ 不会被平移 $\mathbf{a}_2$ 混合，而只是被乘上相位因子
\begin{equation*}
|\mathbf{a}_2\rangle_1=e^{ik_2}|0\rangle_1,\qquad|\mathbf{a}_2\rangle_2=e^{ik_2'}|0\rangle_2.
\tag{1.56}
\end{equation*}
对平移 $\mathbf{a}_3$ 也有类似的结果，于是我们基本上回到了与式 (1.45) 实质相同的情形：对 $|0\rangle_1$ 和 $|0\rangle_2$ 中的每一个都存在一个波矢，使得
\begin{equation*}
|\mathbf{l}\rangle=e^{i\mathbf{k}\cdot\mathbf{l}}|0\rangle.
\tag{1.57}
\end{equation*}
把这一论证推广到两个以上能量简并的态，显然是轻而易举的。

这样，定理便在一般情形下得证，其含义是：波函数的任何一个简并解都可以表示成一些同能量解的线性组合，其中每一个都必须满足这类条件，尽管并非全都取同一个 $\mathbf{k}$ 值。

这就是\emph{Bloch定理}。其实，还存在一个一般的群论证明，它是下述定理的一条推论：“在复数域中，Abel群的任何表示都可以约化为若干一维表示之和”。晶体的平移群是Abel群——也就是说，群的所有元素彼此对易。这来自一个明显的事实：比如平移 $(\mathbf{a}_1+\mathbf{a}_2)$ 与 $(\mathbf{a}_2+\mathbf{a}_1)$ 是同一回事——我们可以沿许多不同的路径到达一个格点，但结果完全一样。矩阵 $\mathsf{T}_1$、$\mathsf{T}_2$ 等等正是这些平移的表示，因而可以约化为对角形式。

\section{约化到一个布里渊区}

Bloch定理的普遍性如此之强，以致在这个阶段我们很难把握它。对电子波而言，它意味着我们可以用波矢 $\mathbf{k}$ 来标记每一个波函数，并写成
\begin{equation*}
\psi_{\mathbf{k}}(\mathbf{r}+\mathbf{l})=e^{i\mathbf{k}\cdot\mathbf{l}}\psi_{\mathbf{k}}(\mathbf{r}).
\tag{1.58}
\end{equation*}
我们注意到，普通的自由电子波满足这个条件：
\begin{equation*}
\psi_{\mathbf{k}}(\mathbf{r})=e^{i\mathbf{k}\cdot\mathbf{r}}.
\tag{1.59}
\end{equation*}
这是意料之中的，因为我们处理的仍然是在周期势中Schrödinger方程的解；只不过这个势碰巧处处为零。空晶格检验（empty-lattice test）在固体理论中往往非常有用。

有时为了方便，人们试图让具有给定 $\mathbf{k}$ 值的电子波函数尽可能看起来像自由电子波。我们令
\begin{equation*}
\psi_{\mathbf{k}}(\mathbf{r})=e^{i\mathbf{k}\cdot\mathbf{r}}u_{\mathbf{k}}(\mathbf{r})
\tag{1.60}
\end{equation*}
并指望 $u_{\mathbf{k}}$ 将近乎恒定。事实上，由式 (1.58) 立刻可以证明，函数 $u_{\mathbf{k}}(\mathbf{r})$ 必须是周期的，即
\begin{equation*}
u_{\mathbf{k}}(\mathbf{r}+\mathbf{l})=u_{\mathbf{k}}(\mathbf{r}).
\tag{1.61}
\end{equation*}
Bloch定理有时就表述成这种形式。

定理中出现的因子 $\exp(i\mathbf{k}\cdot\mathbf{l})$，与我们研究周期函数时遇到的因子 $\exp(i\mathbf{g}\cdot\mathbf{l})$ 相类似。波矢 $\mathbf{k}$ 显然与倒格矢 $\mathbf{g}$ 具有相同的量纲；它属于倒格子空间。如果某个态碰巧具有波矢 $\mathbf{g}$，那它就应该是一个周期函数：
\begin{align*}
\psi_{\mathbf{g}}(\mathbf{r}+\mathbf{l})&=e^{i\mathbf{g}\cdot\mathbf{l}}\psi_{\mathbf{g}}(\mathbf{r})\notag\\
&=\psi_{\mathbf{g}}(\mathbf{r}),
\tag{1.62}
\end{align*}
因为
\begin{equation*}
e^{i\mathbf{g}\cdot\mathbf{l}}=1
\tag{1.63}
\end{equation*}
对一切 $\mathbf{l}$ 成立。

再设态 $\psi_{\mathbf{k}}$ 具有满足下式的波矢 $\mathbf{k}$：
\begin{equation*}
\mathbf{k}=\mathbf{g}+\mathbf{k}',
\tag{1.64}
\end{equation*}
其中 $\mathbf{g}$ 是某个倒格矢，而 $\mathbf{k}'$ 是另一个矢量。于是由式 (1.58) 和 (1.63)，
\begin{align*}
\psi_{\mathbf{k}}(\mathbf{r}+\mathbf{l})&=e^{i(\mathbf{g}+\mathbf{k}')\cdot\mathbf{l}}\psi_{\mathbf{k}}(\mathbf{r})\notag\\
&=e^{i\mathbf{g}\cdot\mathbf{l}}e^{i\mathbf{k}'\cdot\mathbf{l}}\psi_{\mathbf{k}}(\mathbf{r})=e^{i\mathbf{k}'\cdot\mathbf{l}}\psi_{\mathbf{k}}(\mathbf{r}).
\tag{1.65}
\end{align*}
这就是说，态 $\psi_{\mathbf{k}}$ 满足Bloch定理，仿佛它具有波矢 $\mathbf{k}'$ 似的。原来的标记 $\mathbf{k}$ 并不唯一；每一个态都拥有一大群可能的波矢，它们彼此之间相差倒格子的矢量。这当然不与定理矛盾，因为定理只不过断言每个态必须至少有\emph{一个}波矢。

这样我们就面临一个问题：如何唯一地定义一个给定态的波矢？可以拿式 (1.60) 作指导，试着选取 $\mathbf{k}$ 使 $u_{\mathbf{k}}(\mathbf{r})$ 尽可能接近常数——但这是任意的，而且正如我们以后会看到的，甚至会引人误入歧途。正确的做法如下：

考虑在一维情形（参看式 (1.8)–(1.12)）会发生些什么：倒格子的类似物是倒格子长度（reciprocal lattice lengths）的集合
\begin{equation*}
g_n=n\frac{2\pi}{a}.
\tag{1.66}
\end{equation*}
一个态可以被赋予集合
\begin{equation*}
k=n\frac{2\pi}{a}+k',
\tag{1.67}
\end{equation*}
中的任何一个波数；也就是说，$k$ 只是在模 $(2\pi/a)$ 的意义下确定的。图 12 中所有的 $\mathbf{k}$ 点都彼此等价。

\begin{figure}[H]
\centering
\includegraphics[width=0.72\textwidth]{fig_12.pdf}
\caption*{图 12\quad 诸点 $k$ 在一维倒格子中全都约化为 $k'$。}
\end{figure}

很自然地取 $k'$ 作为所有这些点的代表，且使 $|k'|$ 尽可能小。换句话说，我们总是选取落在范围
\begin{equation*}
-\frac{\pi}{a}<k\leqslant\frac{\pi}{a}.
\tag{1.68}
\end{equation*}
之内的值来作为波数。

显然，这个范围对我们的一维系统来说就等价于一个布里渊区；它是倒格子的一个晶胞。在三维情形我们如法炮制——我们选取波矢，使它落在倒格子空间的第一布里渊区内。这件事总是可能的，由初等几何即可明白。我们在式 (1.64) 中这样选取 $\mathbf{g}$ 的值，使 $|\mathbf{k}'|$ 尽可能小——也就是说，使它离倒格子的原点尽可能地近。这意味着 $|\mathbf{k}'|$ 所在之处离原点要比离倒格子的任何其他格点都近——这就等于说它落在该格子的Wigner–Seitz原胞内，也就是落在布里渊区内。显然，我们可以把倒格子空间中的任何一点 $\mathbf{k}$ \emph{约化}为布里渊区中的一点，于是任何态都可以在简约区图式（reduced zone scheme）中获得一个标记。

于是任何态都可以用它的简约波矢（reduced wave-vector）来表征。但可能有许多态具有同一个简约波矢而能量各不相同。因此，采用简约区图式使我们无法给每个态指定一个互不相同的 $\mathbf{k}$ 值。

\begin{figure}[H]
\centering
\includegraphics[width=0.72\textwidth]{fig_13.pdf}
\caption*{图 13\quad 诸波矢 $\mathbf{k}$ 全都约化为位于布里渊区中的 $\mathbf{k}'$。}
\end{figure}

看一看自由电子态在空晶格中会发生些什么，是很有意思的。设
\begin{align*}
\psi&=e^{i\mathbf{k}\cdot\mathbf{r}}\notag\\
&=e^{i(\mathbf{k}-\mathbf{g})\cdot\mathbf{r}}e^{i\mathbf{g}\cdot\mathbf{r}}\notag\\
&=e^{i\mathbf{k}'\cdot\mathbf{r}}\left(e^{i\mathbf{g}\cdot\mathbf{r}}\right),
\tag{1.69}
\end{align*}
其中 $\mathbf{k}'$ 是实际波矢 $\mathbf{k}$ 的约化值。这具有式 (1.60) 的形式，因为 $\exp(i\mathbf{g}\cdot\mathbf{r})$ 是晶格中的周期函数；但它是平面波的一种极其生硬人为的表示。如果在一维情形下能量由
\begin{equation*}
\mathcal{E}(\mathbf{k})=\frac{\hbar^{2}k^{2}}{2m},
\tag{1.70}
\end{equation*}
给出，那么在简约区中它将表现为 $\mathbf{k}'$ 的多值函数。在图 14 中，抛物线的 $AB$ 段经平移一个倒格矢被移到 $A'B'$ 处，如此等等。显然，在这种情形下，其中每个态都用它的“真实”波矢 $\mathbf{k}$ 来表示的扩展区图式（extended zone scheme）要更自然。

\section{边界条件：态的计数}

以上这一切都假定了完全的平移对称性，而这将要求晶格为无限。这在数学上相当麻烦，因为那样一来，原子和有待处理的态都会有无限多个。我们唯一的办法，是先处理一个原子数目有限的系统，然后令这个数目趋于无穷而取极限。我们也许会像对待真实表面那样引入边界，例如要求波函数在边界上必须为零。但这样做同样不便，因为此时系统的严格定态将不得不计及来自边界的所有反射，结果都会成为驻波。传导电子通常在极短的距离内就会被杂质或热振动非相干地散射开；要用这类函数去表示传导电子的态，在数学上笨拙，而且在本质上是不物理的。

\begin{figure}[H]
\centering
\includegraphics[width=0.72\textwidth]{fig_14.pdf}
\caption*{图 14\quad 简约区中自由电子的能量。}
\end{figure}

有一个数学上的手法，能圆满地解决态的计数问题，却又不会引入任何来自边界的直接物理效应。这就是采用循环（cyclic）边界条件，即 Born–von Kármán 边界条件。

在一维情形，设这个“晶体”含有 $L$ 个晶胞，它们首尾相接连成一个环。这一做法对电子波函数的后果，可以用如下条件来表达：
\begin{equation*}
\psi(x+La)=\psi(x),
\tag{1.71}
\end{equation*}
这就保证了在接合点处没有不连续。

但我们的一维 Bloch 条件给出
\begin{equation*}
\psi_k(x+La)=e^{ikLa}\psi_k(x),
\tag{1.72}
\end{equation*}
于是必须有
\begin{equation*}
e^{ikLa}=1,
\tag{1.73}
\end{equation*}
即
\begin{equation*}
k=\frac{2\pi m}{La},
\tag{1.74}
\end{equation*}
其中 $m$ 为整数。

在一维的简约区中，我们有 $-\pi/a<k\leqslant\pi/a$，因此落在范围
\begin{equation*}
-\frac{1}{2}L<m\leqslant\frac{1}{2}L
\tag{1.75}
\end{equation*}
之内的那些整数，将给出所有本质上不同的简约波数值。这样的值恰好有 $L$ 个（这里当然要分别处理 $L$ 为奇数和 $L$ 为偶数的不同情形——一个琐碎的细节），并且相邻两值相隔 $(1/L)(2\pi/a)$。由于假定 $L$ 非常大，这个分布实际上是连续的，并且在倒格子空间中具有恒定的密度。

在三维情形，我们把这一论证加以推广，宣称宏观系统在三个维度上都是循环的。我们把它取为沿晶格三个基矢方向尺寸分别为 $L_1\mathbf{a}_1$、$L_2\mathbf{a}_2$、$L_3\mathbf{a}_3$ 的晶体，并写下
\begin{equation*}
\psi(\mathbf{r}+L_1\mathbf{a}_1)=\psi(\mathbf{r}),\quad
\psi(\mathbf{r}+L_2\mathbf{a}_2)=\psi(\mathbf{r}),\quad
\psi(\mathbf{r}+L_3\mathbf{a}_3)=\psi(\mathbf{r}),
\tag{1.76}
\end{equation*}
正如式 (1.71) 一样。

对于波矢为 $\mathbf{k}$ 的 Bloch 态，这些条件意味着
\begin{equation*}
e^{i\mathbf{k}\cdot(L_1\mathbf{a}_1)}
=e^{i\mathbf{k}\cdot(L_2\mathbf{a}_2)}
=e^{i\mathbf{k}\cdot(L_3\mathbf{a}_3)}=1,
\tag{1.77}
\end{equation*}
而这只有当 $\mathbf{k}$ 取如下形式时方能实现：
\begin{align*}
\mathbf{k} &= k_1\mathbf{b}_1+k_2\mathbf{b}_2+k_3\mathbf{b}_3\\
&= \frac{2\pi m_1}{L_1}\mathbf{b}_1+\frac{2\pi m_2}{L_2}\mathbf{b}_2+\frac{2\pi m_3}{L_3}\mathbf{b}_3,
\tag{1.78}
\end{align*}
其中 $m_1$、$m_2$、$m_3$ 为整数，而矢量 $\mathbf{b}_1$、$\mathbf{b}_2$、$\mathbf{b}_3$ 当然就是 $(\mathbf{a}_1,\mathbf{a}_2,\mathbf{a}_3)$ 的倒格子三矢组（reciprocal triad），如式 (1.19) 中那样。

但这一结果显然可与倒格矢的定义相类比：
\begin{equation*}
\mathbf{g}=n_1\cdot 2\pi\mathbf{b}_1+n_2\cdot 2\pi\mathbf{b}_2+n_3\cdot 2\pi\mathbf{b}_3,
\tag{1.79}
\end{equation*}
其中 $n_1$、$n_2$、$n_3$ 为整数。把倒格子晶胞的生成矢量沿 $\mathbf{b}_1$ 方向分成 $L_1$ 份、沿 $\mathbf{b}_2$ 方向分成 $L_2$ 份、沿 $\mathbf{b}_3$ 方向分成 $L_3$ 份，便得到式 (1.78) 中 $\mathbf{k}$ 的允许值（allowed values）。于是，宛如生出一阵细密的疹点，点子均匀地布满倒格子空间，如图 15 所示。

要计算这些点的密度，只需注意到：让整数 $m_1$、$m_2$、$m_3$ 取遍
\begin{equation*}
0\leqslant m_1<L_1,\quad
0\leqslant m_2<L_2,\quad
0\leqslant m_3<L_3.
\tag{1.80}
\end{equation*}
即可铺满倒格子的一个晶胞。

但这并不是我们自然会选来作为 $\mathbf{k}$ 矢量基本区的晶胞。或许改取如下范围更好：
\begin{equation*}
-\frac{1}{2}L_1<m_1<\frac{1}{2}L_1,\quad
-\frac{1}{2}L_2<m_2<\frac{1}{2}L_2,\quad
-\frac{1}{2}L_3<m_3<\frac{1}{2}L_3,
\tag{1.81}
\end{equation*}
以与式 (1.75) 类比。这将给出一个以原点为中心的平行六面体晶胞。然而，晶胞的最佳选择是倒格子的 Wigner–Seitz 原胞，也就是我们的老朋友——布里渊区。

\begin{figure}[H]
\centering
\includegraphics[width=0.72\textwidth]{fig_15.pdf}
\caption*{图 15\quad 布里渊区与倒格子的平行六面体晶胞覆盖相同的面积，且恰含 N 个“允许的 k 矢量”。}
\end{figure}

由于布里渊区的体积与平行六面体晶胞相同，它必定恰好含有同样数目的 $\mathbf{k}$ 的“允许值”。由式 (1.80) 或式 (1.81)，这个数目显然就是 $L_1\times L_2\times L_3=N$，而这正是我们整块宏观晶体中晶胞的数目。这是一个极其重要的定理：\emph{一块晶体中有多少个晶胞，一个布里渊区中就恰有多少个允许的波矢。}

我们已经证明过（式 (1.32)），一个区的体积为 $8\pi^3/v_c$。若体积为 $V$ 的晶体中有 $N$ 个晶胞，则
\begin{equation*}
v_c=V/N.
\tag{1.82}
\end{equation*}
于是，每个允许的 $\mathbf{k}$ 矢量在 $\mathbf{k}$ 空间中占据的体积便是
\begin{equation*}
\frac{1}{N}\,\frac{8\pi^3}{v_c}=\frac{8\pi^3}{V}.
\tag{1.83}
\end{equation*}
因此，\emph{倒格子空间单位体积内有 $V/8\pi^3$ 个允许的 $\mathbf{k}$ 矢量。}

实际上 $N$ 非常大，以致这一分布被当作连续分布来处理。我们常把对 $\mathbf{k}$ 矢量的求和写成一个积分
\begin{equation*}
\sum_{\mathbf{k}}\rightarrow\int d\mathbf{k}\equiv\frac{V}{8\pi^3}\iiint d^3k,
\tag{1.84}
\end{equation*}
这里用单个积分作为求和之极限的简明记号。然而，重要的是要记住：当我们真正要在 $\mathbf{k}$ 空间中做具体求积时，必须计入权重因子 $V/8\pi^3$。为简单起见，我们通常将假定 $V=1$，这时 $N$ 就是单位体积晶体中的晶胞数，而 $1/N$ 就是一个晶胞的体积。

结果 (1.83) 实际上不依赖于任何所假定的区结构；它在辐射理论中以及在自由电子情形下早已为人熟知。但就许多目的而言，布里渊区恰好含有 $N$ 个允许点这一事实更为有用。它表明这个区在本质上是不变的，仅依赖于晶体结构；增大整块晶体的尺寸，只会使 k 空间中的态密度（density of states）增大。

不过，必须承认，Born–von Kármán 条件在物理上是无法实现的。在二维情形，画在环面（torus）表面上的晶胞网络沿两个方向都是循环的；但对三维晶格来说，无论怎样施展拓扑上的扭曲，都无法使三个条件同时满足。尽管如此，它作为一种数学技巧用起来极为出色，而且可以用一些繁冗的数学来为它提供证明。事实是：在大量子数的渐近极限下，态密度对边界条件的具体形式非常不敏感。只要我们并不打算讨论边界效应本身，循环边界条件就远远是最方便的选择。
